\documentclass{article}
\usepackage[utf8]{inputenc}
\usepackage{amsmath}
\usepackage{geometry}
\geometry{margin=2cm}
\begin{document}

\section*{Questão 105 — ENEM 2024 — Ciências da Natureza}

\subsection*{Temática: Vírus, linfócitos T e câncer \emoji{🦠}}

\paragraph{Enunciado resumido:} O HTLV-1 estimula a produção excessiva de linfócitos T. Qual a possível consequência desse efeito?

\section*{1. Interpretação e Estratégia de Resolução}

A questão pede a compreensão da relação entre o efeito fisiopatológico descrito (proliferação de linfócitos T induzida pelo HTLV-1) e a consequência patológica mais provável. A estratégia consiste em:
\begin{itemize}
  \item Relacionar o mecanismo (proliferação de linfócitos T) ao tipo de doença esperada (neoplasias hematológicas).
  \item Eliminar alternativas que não se correlacionam causalmente com o efeito recebido.
  \item Saber distinguir entre vírus retrovirais, associando cada um ao respectivo quadro clínico típico.
\end{itemize}

\section*{2. Conteúdo e Mini-revisão}

\begin{itemize}
  \item \textbf{HTLV-1:} Retrovírus que infecta linfócitos T, estimulando sua proliferação.
  \item \textbf{Câncer hematológico:} Neoplasias como linfoma e leucemia surgem pela multiplicação descontrolada dessas células.
  \item \textbf{Diferença-chave:} HIV destrói linfócitos T (causa AIDS), enquanto o HTLV-1 estimula crescimento dessas células (pode causar câncer).
\end{itemize}

\begin{center}
\begin{tabular}{|l|l|l|}
\hline
\textbf{Vírus} & \textbf{Célula-alvo} & \textbf{Efeito} \\
\hline
HIV  & Linfócito T      & Destruição (imunodeficiência) \\
HTLV-1 & Linfócito T      & Proliferação (câncer hematológico) \\
\hline
\end{tabular}
\end{center}

\section*{3. Resolução Detalhada}

1. Identifique o mecanismo: HTLV-1 aumenta linfócitos T.
2. Proliferação celular descontrolada leva a câncer.
3. Descartando as demais opções (AIDS, diabetes, hepatite B, hemorragia) que não se relacionam à proliferação celular dos linfócitos T.
4. Assim, a alternativa B) câncer é a única correta.

\textbf{Pulo do gato:} Sempre que um agente infeccioso induzir proliferação celular desordenada, pense em neoplasia como consequência principal.

\section*{4. Armadilhas Comuns na Questão}

\begin{itemize}
  \item \textbf{Confundir com HIV:} HTLV-1 não causa imunossupressão, mas sim proliferação anormal.
  \item \textbf{Marcar doenças conhecidas:} Como diabetes e hepatite B, pela familiaridade do nome e não pelo mecanismo envolvido.
  \item \textbf{Associar distúrbio hematológico à hemorragia:} Nem todo distúrbio do sangue leva à hemorragia, especialmente quando há proliferação celular sem destruição.
\end{itemize}

\section*{5. Código Oficial da Questão}

\texttt{CIENCIASNATURIZA\_VIROLOGIA\_CLINICA\_001}

\section*{6. Alternativas e Explicações}
\begin{description}
  \item[A)] \textbf{AIDS:} Errada. HIV destrói linfócitos T e causa imunodeficiência. HTLV-1 estimula proliferação.
  \item[B)] \textbf{Câncer:} Correta. HTLV-1 pode causar cânceres hematológicos como leucemia e linfoma de células T.
  \item[C)] \textbf{Diabetes:} Errada. Não há relação causal entre linfócitos T e diabetes.
  \item[D)] \textbf{Hepatite B:} Errada. O HTLV-1 não se relaciona a doenças hepáticas.
  \item[E)] \textbf{Hemorragia:} Errada. O aumento de linfócitos T não provoca sangramento, mas sim risco de câncer.
\end{description}

\section*{Resumo para o professor}
Esta questão avalia associação entre mecanismo fisiopatológico (proliferação de linfócitos) e consequência clínica (câncer hematológico). Prefira sempre análise causal ao escolher entre doenças comuns.

\end{document}